\begin{proof}[Proof of \cref{obj:lem:legal-cancellation}]
\begin{enumerate}
\item Define the shifted polynomials and the untapered normalized block by
\[
R_j(t)=P_j^{\mathrm{Leg}}(2t-1),
\qquad
q_m^{\mathrm{pol}}(t)
=\frac1{N_m}\sum_{j=m}^{2m}(2j+1)(-1)^jR_j(t),
\qquad j\in\mathbb N,\quad t\in\mathbb R.
\]
Thus \(g_m(t)=(1-t)q_m^{\mathrm{pol}}(t)\). Summing the odd weights gives
\[
\sum_{j=m}^{2m}(2j+1)
=\sum_{j=0}^{2m}(2j+1)-\sum_{j=0}^{m-1}(2j+1)
=(2m+1)^2-m^2
=N_m=3m^2+4m+1>0.
\]
The endpoint identity in \cref{obj:lem:classical-legendre-facts} therefore yields
\[
q_m^{\mathrm{pol}}(0)
=\frac1{N_m}\sum_{j=m}^{2m}(2j+1)(-1)^j(-1)^j
=1.
\]
Consequently,
\[
g_m(0)=1,\qquad g_m(1)=0.
\]
% lean: CausalSmith.Stat.RdTruesideNoiseFrontier.block_zero
% lean: CausalSmith.Stat.RdTruesideNoiseFrontier.block_one

\item For \(t\in[0,1]\), the argument \(2t-1\) lies in \([-1,1]\). The unit bound in \cref{obj:lem:classical-legendre-facts} gives
\[
\left|\sum_{j=m}^{2m}(2j+1)(-1)^jR_j(t)\right|
\le\sum_{j=m}^{2m}(2j+1)|R_j(t)|
\le\sum_{j=m}^{2m}(2j+1)=N_m.
\]
Since \(N_m>0\) and \(0\le1-t\le1\), it follows that
\[
|q_m^{\mathrm{pol}}(t)|\le1,
\qquad
|g_m(t)|=(1-t)|q_m^{\mathrm{pol}}(t)|\le1.
\]
% lean: CausalSmith.Stat.RdTruesideNoiseFrontier.block_abs_le_one

\item Applying the substitution \(u=2t-1\) to the orthogonality identity in \cref{obj:lem:classical-legendre-facts} gives, for \(j,k\in\mathbb N\),
\[
\int_0^1R_j(t)R_k(t)\,dt
=
\begin{cases}
(2j+1)^{-1},&j=k,\\
0,&j\ne k.
\end{cases}
\]
Expanding the square of the finite sum, the off-diagonal terms vanish and each diagonal term contributes \(2j+1\). Hence
\[
\begin{aligned}
\int_0^1[q_m^{\mathrm{pol}}(t)]^2\,dt
&=\frac1{N_m^2}
\sum_{j=m}^{2m}\sum_{k=m}^{2m}
(2j+1)(2k+1)(-1)^{j+k}
\int_0^1R_j(t)R_k(t)\,dt\\
&=\frac1{N_m^2}\sum_{j=m}^{2m}(2j+1)
=\frac1{N_m}.
\end{aligned}
\]
All these integrands are polynomials, so finite sums and integrals may be interchanged. Finally,
\[
g_m(t)^2=(1-t)^2[q_m^{\mathrm{pol}}(t)]^2
\le[q_m^{\mathrm{pol}}(t)]^2
\qquad(t\in[0,1]),
\]
and integration proves
\[
\int_0^1g_m(t)^2\,dt\le\frac1{N_m}.
\]
% lean: CausalSmith.Stat.RdTruesideNoiseFrontier.block_sq_integral_le

\item Fix an integer \(k\) with \(0\le k\le m-2\), and define
\[
\varphi_k(t)=t^k(1-t),\qquad t\in\mathbb R.
\]
Its degree is at most \(k+1\). For every \(j\in\{m,\ldots,2m\}\),
\[
\deg\varphi_k\le k+1<m\le j.
\]
The reversed shifted polynomials \((-1)^iR_i\), \(0\le i<j\), have respective degrees \(i\) and nonzero leading coefficients. Successively cancelling the leading term of a polynomial of degree less than \(j\) therefore expresses it as a linear combination of these polynomials; the zero polynomial is represented by zero coefficients. In particular, there exist real coefficients \(c_{k,j,i}\), \(0\le i<j\), such that
\[
\varphi_k(t)=\sum_{i=0}^{j-1}c_{k,j,i}(-1)^iR_i(t)
\qquad(t\in\mathbb R).
\]
The shifted orthogonality identity established above now gives
\[
\int_0^1\varphi_k(t)R_j(t)\,dt
=\sum_{i=0}^{j-1}c_{k,j,i}(-1)^i
\int_0^1R_i(t)R_j(t)\,dt
=0.
\]
Using the defining finite sum for \(g_m\), we conclude that
\[
\int_0^1t^kg_m(t)\,dt
=\frac1{N_m}\sum_{j=m}^{2m}(2j+1)(-1)^j
\int_0^1\varphi_k(t)R_j(t)\,dt
=0.
\]
This includes \(m=2\), for which the permitted moment is \(k=0\).
% lean: CausalSmith.Stat.RdTruesideNoiseFrontier.block_moment_zero

\item Differentiating the finite polynomial sum gives
\[
(q_m^{\mathrm{pol}})'(t)
=\frac1{N_m}\sum_{j=m}^{2m}(2j+1)(-1)^jR_j'(t).
\]
Every index in this sum satisfies \(j\ge m\ge2\), so \cref{obj:lem:shifted-legendre-derivative} supplies
\[
|R_j'(t)|\le j(j+1)\le(2m)(2m+1)
\qquad(t\in[0,1]).
\]
Thus
\[
\begin{aligned}
|(q_m^{\mathrm{pol}})'(t)|
&\le\frac1{N_m}\sum_{j=m}^{2m}(2j+1)|R_j'(t)|\\
&\le\frac1{N_m}\sum_{j=m}^{2m}(2j+1)(2m)(2m+1)
=(2m)(2m+1).
\end{aligned}
\]
The product rule yields
\[
g_m'(t)=-q_m^{\mathrm{pol}}(t)
+(1-t)(q_m^{\mathrm{pol}})'(t).
\]
Using the unit bound and \(0\le1-t\le1\), we obtain
\[
|g_m'(t)|
\le1+(2m)(2m+1)
=1+4m^2+2m
\le7m^2.
\]
The final calibration follows from
\[
7m^2-(1+4m^2+2m)=(3m+1)(m-1)\ge0
\qquad(m\ge2).
\]
% lean: CausalSmith.Stat.RdTruesideNoiseFrontier.block_deriv_bound

\item Define the clamped, rescaled argument by
\[
\mathfrak c_b(x)=\max\{0,\min\{1,x/b\}\},
\qquad x\in\mathbb R.
\]
It belongs to \([0,1]\). For \(x<0\) it equals zero, for \(0\le x\le b\) it equals \(x/b\), and for \(x>b\) it equals one. The endpoint identities and \cref{obj:def:legal-extension} therefore give
\[
v_{b,m}(x)=g_m(\mathfrak c_b(x)),
\qquad
|v_{b,m}(x)|\le1
\qquad(x\in[-1,1]).
\]
Both minimum and maximum with a fixed real number are contractions. More explicitly, for \(x,t\in\mathbb R\),
\[
\begin{aligned}
|\mathfrak c_b(x)-\mathfrak c_b(t)|
&\le|\min\{1,x/b\}-\min\{1,t/b\}|\\
&\le|x/b-t/b|
=\frac{|x-t|}{b}.
\end{aligned}
\]
The mean value bound for the polynomial \(g_m\) on the convex interval \([0,1]\) gives
\[
|g_m(t_1)-g_m(t_2)|
\le7m^2|t_1-t_2|
\qquad(t_1,t_2\in[0,1]).
\]
Combining these two inequalities proves the bound across all pieces of the extension:
\[
|v_{b,m}(x)-v_{b,m}(t)|
\le\frac{7m^2}{b}|x-t|
\qquad(x,t\in[-1,1]).
\]

For the Hölder estimate, define
\[
\Lambda_{b,m}=\frac{m^2}{b}>0,
\qquad
\Lambda_{b,m}^{\beta}=(b/m^2)^{-\beta}.
\]
Fix distinct \(x,t\in[-1,1]\). If \(\Lambda_{b,m}|x-t|\le1\), then \(0<\beta\le1\) implies
\[
\Lambda_{b,m}|x-t|
\le(\Lambda_{b,m}|x-t|)^\beta,
\]
so the Lipschitz estimate gives
\[
|v_{b,m}(x)-v_{b,m}(t)|
\le7\Lambda_{b,m}|x-t|
\le7(\Lambda_{b,m}|x-t|)^\beta.
\]
If \(\Lambda_{b,m}|x-t|>1\), then
\[
1\le(\Lambda_{b,m}|x-t|)^\beta,
\qquad
|v_{b,m}(x)-v_{b,m}(t)|
\le|v_{b,m}(x)|+|v_{b,m}(t)|\le2.
\]
Consequently, in this case as well,
\[
|v_{b,m}(x)-v_{b,m}(t)|
\le2\le7(\Lambda_{b,m}|x-t|)^\beta.
\]
Since the factors are positive,
\[
(\Lambda_{b,m}|x-t|)^\beta
=\Lambda_{b,m}^\beta|x-t|^\beta.
\]
Dividing by \(|x-t|^\beta\) and taking the supremum proves
\[
[v_{b,m}]_\beta\le7(b/m^2)^{-\beta}.
\]
The same argument applies at \(\beta=1\).
% lean: CausalSmith.Stat.RdTruesideNoiseFrontier.legalExtension_holder_bound

\item For the probability profiles in \cref{obj:def:alternatives}, write
\[
p_\pm(x)=\frac12\pm\kappa(b/m^2)^\beta v_{b,m}(x),
\qquad x\in[-1,1].
\]
The clamp representation makes \(v_{b,m}\), and hence \(p_\pm\), continuous and Borel measurable on this interval. Since \(m\ge2\) and \(0<b\le1\),
\[
0<\frac b{m^2}\le1,
\qquad
0<(b/m^2)^\beta\le1.
\]
Together with the unit bound, this implies
\[
|(b/m^2)^\beta v_{b,m}(x)|\le1,
\qquad
\frac12-\frac1{100}\le p_\pm(x)\le\frac12+\frac1{100}.
\]
In particular,
\[
\frac14\le p_\pm(x)\le\frac34,
\]
so both profiles define valid Bernoulli probabilities.

The construction in \cref{obj:def:alternatives} uses an independent Gaussian factor. Its Gaussian marginal and product structure give
\[
\varepsilon\sim N(0,1),
\qquad
\varepsilon\perp(X,Y^0,Y^1),
\]
as required by \cref{obj:ass:gaussian,obj:ass:error-independence}. The uniform score and binary Bernoulli outcomes give the latent support in \cref{obj:def:model}. Moreover,
\[
f(P^\pm_{b,m},x)=\frac12
\qquad(x\in[-1,1]),
\qquad
\int_{-1}^1f(P^\pm_{b,m},x)\,dx=1.
\]
This density is continuous on the latent interval and satisfies \cref{obj:ass:density-bounds}. The conditional Bernoulli means are
\[
\mu_0(P^\pm_{b,m},x)=\frac12,
\qquad
\mu_1(P^\pm_{b,m},x)=p_\pm(x),
\]
which are continuous and satisfy \cref{obj:ass:mean-bounds}.

For distinct \(x,t\in[-1,1]\), the preceding Hölder estimate gives
\[
\begin{aligned}
|p_\pm(x)-p_\pm(t)|
&=\kappa(b/m^2)^\beta
|v_{b,m}(x)-v_{b,m}(t)|\\
&\le7\kappa(b/m^2)^\beta(b/m^2)^{-\beta}|x-t|^\beta\\
&=\frac7{100}|x-t|^\beta
\le2|x-t|^\beta,
\end{aligned}
\]
where
\[
(b/m^2)^\beta(b/m^2)^{-\beta}=1
\]
because \(b/m^2>0\). When \(x=t\), the same inequality follows from the zero difference. The untreated mean is constant, so
\[
[\mu_0(P^\pm_{b,m},\cdot)]_\beta=0,
\qquad
[\mu_1(P^\pm_{b,m},\cdot)]_\beta\le\frac7{100}\le2.
\]
Thus \cref{obj:ass:mean-holder} holds, completing all conditions of \cref{obj:def:model} for each sign and every \(\sigma\in\mathbb R\):
\[
P^\pm_{b,m}\in\mathcal M_\beta(\sigma).
\]
% lean: CausalSmith.Stat.RdTruesideNoiseFrontier.altLaw_model

\item At the cutoff,
\[
v_{b,m}(0)=g_m(0)=1.
\]
Evaluating the conditional means therefore gives
\[
\theta(P^\pm_{b,m})
=\mu_1(P^\pm_{b,m},0)-\mu_0(P^\pm_{b,m},0)
=\pm\kappa(b/m^2)^\beta.
\]
Since \(\kappa(b/m^2)^\beta\ge0\), subtraction and the absolute value yield
\[
\left|\theta(P^+_{b,m})-\theta(P^-_{b,m})\right|
=\left|2\kappa(b/m^2)^\beta\right|
=2\kappa(b/m^2)^\beta.
\]
% lean: CausalSmith.Stat.RdTruesideNoiseFrontier.altLaw_separation
\end{enumerate}
\end{proof}
